%% ============================================================
%% Negative-trace construction
%% ============================================================
\begin{figure}[H]
\centering
% Match the overview and data figures without changing the document palette.
\definecolor{paperink}{HTML}{1F3128}
\definecolor{codegray}{HTML}{66786F}
\definecolor{codegreen}{HTML}{2E7D5B}
\definecolor{codered}{HTML}{B85C47}
\definecolor{wrongred}{HTML}{D99382}
\definecolor{promptframe}{HTML}{D8E2DC}
\definecolor{algobg}{HTML}{F2F6F3}
\renewcommand{\tiny}{\fontsize{7.2}{8.4}\selectfont}
\renewcommand{\scriptsize}{\fontsize{8.6}{9.8}\selectfont}
\resizebox{\textwidth}{!}{%
\begin{tikzpicture}[
  font=\sffamily\footnotesize,
  >=Stealth,
  flow/.style={-{Stealth[length=2.2mm]}, draw=paperink!58, line width=1.1pt},
  panel/.style={rounded corners=3pt, fill=algobg, draw=promptframe, line width=0.55pt},
  subcard/.style={rounded corners=3.5pt, fill=white, draw=none},
  stage/.style={font=\sffamily\scriptsize\bfseries, text=paperink},
  ttl/.style={font=\sffamily\scriptsize\bfseries, text=paperink, anchor=west},
  ann/.style={font=\sffamily\tiny, text=codegray, align=center},
  pdot/.style={circle, fill=codegreen, inner sep=1.55pt},
  neg/.style={circle, draw=codered, fill=wrongred!35, line width=0.8pt,
    inner sep=1.35pt},
  pedge/.style={-{Stealth[length=1.5mm]}, draw=codegreen!75, line width=0.75pt},
  nedge/.style={-{Stealth[length=1.5mm]}, draw=codered!80, line width=0.75pt},
  perturb/.style={-{Stealth[length=1.6mm]}, draw=codered!80, line width=0.8pt,
    dashed}
]

%% ---------------- Stage 1: execute ----------------
\node[stage] at (1.30,1.62) {Execute};
\draw[panel] (0,-1.30) rectangle (2.60,1.30);
\node[ann, text=paperink] at (1.30,0.88)
  {masked loop\\[-1pt]{\color{codered}\(\mathcal G\) hidden}};
\node[pdot] (e1) at (0.50,0.02) {};
\node[pdot] (e2) at (1.15,0.02) {};
\node[pdot] (e3) at (1.80,0.02) {};
\node[pdot] (e4) at (2.25,0.02) {};
\draw[pedge] (e1) -- (e2);
\draw[pedge] (e2) -- (e3);
\draw[pedge] (e3) -- (e4);
\node[ann] at (1.30,-0.78) {observed reachable\\[-1pt]states \(\mathcal E\)};

\draw[flow] (2.60,0) -- (3.30,0);

%% ---------------- Stage 2: two families ----------------
\node[stage] at (6.65,2.18) {Perturb};
\draw[panel] (3.30,-1.90) rectangle (10.00,1.90);

% relation family
\draw[subcard] (3.55,0.12) rectangle (9.75,1.66);
\node[ttl, text=paperink] at (3.75,1.42) {Relation};
\node[pdot] (r1) at (4.15,0.62) {};
\node[pdot] (r2) at (5.05,0.62) {};
\node[pdot] (r3) at (5.95,0.62) {};
\draw[pedge] (r1) -- (r2);
\draw[pedge] (r2) -- (r3);
\node[ann] at (5.05,0.36) {\((2,3)\)};
\node[neg] (roff) at (6.75,1.18) {};
\draw[perturb] (r2) -- (roff);
\node[ann, text=codered] at (7.32,1.18) {\((3,3)\)};
\node[ann, anchor=east] at (9.62,0.70)
  {guard preserved,\\[-1pt]off the transition\\[-1pt]tangent};

% post-exit family
\draw[subcard] (3.55,-1.66) rectangle (9.75,-0.12);
\node[ttl, text=paperink] at (3.75,-0.38) {Post-exit};
\node[pdot] (p1) at (4.15,-1.10) {};
\node[pdot] (p2) at (4.95,-1.10) {};
\node[pdot] (p3) at (5.75,-1.10) {};
\draw[pedge] (p1) -- (p2);
\draw[pedge] (p2) -- (p3);
\draw[dashed, draw=codegray!85, line width=0.7pt] (6.15,-1.50) -- (6.15,-0.62);
\node[ann] at (6.15,-0.42) {exit};
\node[neg] (q1) at (6.60,-1.10) {};
\node[neg] (q2) at (7.40,-1.10) {};
\node[neg] (q3) at (8.20,-1.10) {};
\draw[nedge] (p3) -- (q1);
\draw[nedge] (q1) -- (q2);
\draw[nedge] (q2) -- (q3);
\node[ann, anchor=east] at (9.62,-1.08) {real body\\[-1pt]continues};

\draw[flow] (10.00,0) -- (10.70,0);

%% ---------------- Stage 3: clean + budget ----------------
\node[stage] at (12.10,1.62) {Clean \& budget};
\draw[panel] (10.70,-1.30) rectangle (13.50,1.30);
\node[ann, anchor=west] at (10.95,0.82)
  {{\color{codered}\(\times\)} observed state};
\node[ann, anchor=west] at (10.95,0.44)
  {{\color{codered}\(\times\)} fresh loop entry};
\node[ann, anchor=west] at (10.95,0.06)
  {{\color{codered}\(\times\)} duplicate};
\draw[draw=promptframe, line width=0.6pt] (10.95,-0.28) -- (13.25,-0.28);
\node[ann, text=paperink] at (12.10,-0.62)
  {48 relation \(\cdot\) 12 post-exit};
\node[ann] at (12.10,-0.98) {nearest perturbations first};

\draw[flow] (13.50,0) -- (14.20,0);

%% ---------------- Output ----------------
\node[circle, fill=wrongred!25,
  inner sep=0pt, minimum size=1.18cm] (out) at (14.95,0)
  {\(\mathcal N\)};
\node[ann] at (14.95,-0.98) {negative traces\\[-1pt]for the RL reward};

\end{tikzpicture}
}%
\caption{\textbf{Construction of target-hidden negative traces.}
Concrete executions provide observed reachable states.  Two transformations
probe relational structure and behavior past a genuine exit.  Conservative
cleaning and per-family budgets produce \(\mathcal N\), which is used only
for RL reward computation and never exposes \(\mathcal G\).}
\label{fig:negative-sampling}
\end{figure}
